% Типовой расчёт № 1 по теории вероятностей
% Лаборатория математики ФН1
% Компиляция: pdflatex tr-01-usloviya.tex

\documentclass[12pt, a4paper]{article}

\usepackage[T2A]{fontenc}
\usepackage[utf8]{inputenc}
\usepackage[russian]{babel}
\usepackage[left=25mm, right=20mm, top=20mm, bottom=20mm]{geometry}
\usepackage{amsmath, amssymb}
\usepackage{enumitem}

\pagestyle{plain}
\setlength{\parindent}{1.25cm}

\begin{document}

\begin{center}
  {\large\bfseries Типовой расчёт № 1}

  \vspace{0.3em}
  {\large Случайные события и вероятность}

  \vspace{0.5em}
  {\small Теория вероятностей и математическая статистика · 2 курс, 4 семестр\\
  Лаборатория математики ФН1}
\end{center}

\vspace{1em}

\noindent
Номер варианта совпадает с порядковым номером студента в журнале группы.
Числовые данные вариантов — в файле \texttt{tr-01-varianty.csv}.

\section*{Задача 1. Классическое определение вероятности}

В урне $n$ шаров, из них $m$ белых. Наудачу извлекают $k$ шаров.
Найти вероятность того, что среди извлечённых окажется ровно два белых шара.
Решение сопроводить схемой подсчёта числа исходов.

\section*{Задача 2. Теоремы сложения и умножения}

Три стрелка независимо стреляют по мишени. Вероятности попадания равны
$p_1$, $p_2$ и $p_3$. Найти вероятность того, что:
\begin{enumerate}[label={\asbuk{enumi})}]
  \item в мишень попадёт ровно один стрелок;
  \item в мишень попадёт хотя бы один стрелок;
  \item в мишень попадут все трое.
\end{enumerate}

\section*{Задача 3. Формула полной вероятности}

Изделия поступают с трёх линий в пропорции $a : b : c$. Доли брака
на линиях равны $q_1$, $q_2$ и $q_3$. Наудачу взятое изделие оказалось
бракованным. Найти вероятность того, что оно изготовлено на второй линии
(формула Байеса).

\section*{Задача 4. Схема Бернулли}

Проводится $N$ независимых испытаний с вероятностью успеха $p$ в каждом.
Требуется:
\begin{enumerate}[label={\asbuk{enumi})}]
  \item найти вероятность ровно $r$ успехов;
  \item найти наивероятнейшее число успехов;
  \item оценить вероятность того, что число успехов отклонится от среднего
        не более чем на 2, воспользовавшись локальной теоремой Муавра~---~Лапласа.
\end{enumerate}

\section*{Задача 5. Дискретная случайная величина}

Случайная величина $X$ задана рядом распределения из файла вариантов.
Построить функцию распределения $F(x)$, изобразить её график,
вычислить $M[X]$, $D[X]$ и $\sigma[X]$.

\section*{Задача 6. Непрерывная случайная величина}

Плотность распределения случайной величины имеет вид
\[
  f(x) =
  \begin{cases}
    0, & x \le 0,\\
    \lambda x e^{-\lambda x}, & x > 0.
  \end{cases}
\]
Найти функцию распределения, математическое ожидание, дисперсию
и вероятность $P(X > M[X])$.

\vspace{1.5em}

\section*{Требования к оформлению}

\begin{itemize}[nosep]
  \item условие каждой задачи переписывается полностью с числами своего варианта;
  \item решение содержит обоснование, а не только вычисления;
  \item ответ выделяется отдельной строкой;
  \item графики строятся на координатной сетке с указанием масштаба.
\end{itemize}

\section*{Критерии оценивания}

Каждая задача оценивается в 2 балла: 2~---~решена верно с обоснованием,
1~---~верный метод при вычислительной ошибке, 0~---~неверный метод либо
решение отсутствует. Расчёт зачтён при 8 баллах из 12 и успешной защите.

\vspace{0.7em}
\noindent\textbf{Сроки.} Выдача~---~3-я учебная неделя,
сдача~---~до конца 9-й недели.

\end{document}
